% Sections 8.1 to 8.3, from lectures/L08/appendix/a_state_machines.md, parts A.1 to A.3: the
% machine designed by hand.

\section{What Is a Finite State Machine}
\label{c8:a1}
A finite state machine (FSM) is a circuit whose behaviour depends not only on its current inputs but
on a finite amount of \emph{history}, summarized in a single value called its \textbf{state}. At any
moment the machine is in exactly one state; on each clock edge it may move to a new one, decided by
its current state and inputs; and its output is computed from that state.

None of this is new hardware. An FSM combines two things you already know: a \textbf{register}
(Chapter~\ref{c3:ch}) holding the current state, one flip-flop per state bit, and
\textbf{combinational logic} computing the next state and the output from the current state and the
inputs.

What is new is the discipline. Instead of ad hoc logic feeding an ad hoc set of flip-flops, you work
in a fixed order: first draw a \textbf{state diagram}, a few named states with labelled transitions
between them, and only once that is right derive the state table, then the next-state logic, then
the gate network or VHDL.

Every synchronous circuit so far is technically a trivial state machine: a counter's state is its
count, a shift register's is its bit pattern. What sets an FSM apart is that the state is not a
count or a shifted pattern but an arbitrary label (\code{STATE_OFF}, \code{STATE_BLINK}, \dots)
whose meaning you choose and whose next value can depend on the input in any way you like.

Two flavours exist, and the distinction decides what an output may look at:
\begin{itemize}
  \item \textbf{Moore machines}: the output depends \emph{only} on the current state. That is the
    kind designed by hand here.
  \item \textbf{Mealy machines}: the output depends on the current state \emph{and} the current
    input. \Secref{c8:a5} covers these briefly, in VHDL only, so you recognize one when you meet it
    and know the one-cycle difference.
\end{itemize}

\section[Hand-Designing a Moore Machine]{Hand-Designing a Moore Machine: State Diagram, State
Table, and Karnaugh-Derived Logic}
\label{c8:a2}
The same process as every gate network in Chapters~\ref{c1:ch} and~\ref{c2:ch}, applied to a
circuit with memory. The worked example is a small, closed, four-state Moore machine: a circular
counter through a Gray-code sequence, advanced one step per button press.

\subsection{Specification}
\label{c8:spec}
\begin{itemize}
  \item Four states, each represented by two flip-flops \code{\{Q1, Q2\}}:

    \begin{narrowtable}{@{}ll@{}}
      \thead{State} & \thead{\code{\{Q1, Q2\}}} \\
      \midrule
      \code{STATE_0} & \code{00} \\
      \code{STATE_1} & \code{01} \\
      \code{STATE_2} & \code{11} \\
      \code{STATE_3} & \code{10} \\
    \end{narrowtable}
  \item A system clock \code{clock}, and an active-low reset \code{reset_n} returning the machine to
    \code{STATE_0}.
  \item An active-low push button \code{button_n}, from which a synchronized, edge-detected pulse
    \code{X} is derived exactly the way Chapter~\ref{c4:ch}'s double-flop synchronizer does.
    \code{X} is high for exactly one clock cycle on a falling edge of \code{button_n}.
  \item An output \code{Y}, connected to an LED, high only in \code{STATE_3}.
  \item The machine is \textbf{closed}: the state after \code{STATE_3} is \code{STATE_0}.
  \item It is a \textbf{Moore} machine: \code{Y} depends only on \code{\{Q1, Q2\}}, never on
    \code{X}.
\end{itemize}

\noindent The encoding is a Gray code (\code{00 -> 01 -> 11 -> 10 -> 00}), where consecutive states
differ in exactly one bit. That is a deliberate choice, not a coincidence:
\begin{itemize}
  \item With an ordinary binary count (\code{00 -> 01 -> 10 -> 11}), a transition like
    \code{01 -> 10} flips two bits at once, and the two flip-flops will not switch at
    \emph{exactly} the same instant.
  \item That skew is harmless to the machine itself. Both flip-flops update on the same clock edge,
    and nothing samples the state register in between, so the machine never acts on the
    intermediate value. This is the same argument as the critical path of \secref{c5:a3}: as long
    as everything settles before the next edge, what happens during the settling does not matter.
  \item What the skew \emph{can} disturb is anything reading the state combinationally. An output
    decoded straight from the state bits, like \code{Y} below, can glitch briefly while they skew.
    So can a value sampled by a \emph{different} clock domain, which is the case \secref{c4:a3}
    warned about.
  \item A Gray code removes both, since only one bit ever changes per transition. That is why it
    matters most on a counter crossing clock domains, and least on a machine whose output is
    registered.
\end{itemize}
Encoding is a design decision in general, and this is the one place in the book you make it by
hand. Binary uses the fewest flip-flops; Gray minimizes bits changing per transition;
\textbf{one-hot} spends one flip-flop per state so the next-state logic is a single wide OR per
state, which is often fastest on an FPGA where flip-flops are plentiful. Once you write the machine
in VHDL as an enumerated type (\secref{c8:a4}), you hand this decision to the synthesis tool, which
picks an encoding for you and re-encodes freely unless you tell it otherwise.

\subsection{State diagram}
\label{c8:diagram}
\bookimage{0.62}{lectures/L08/appendix/images/fsm_state_diagram.png}{The four-state Gray-code Moore
machine: each state holds on \code{X = 0} and advances on \code{X = 1}, and each bubble reads
\code{Q1Q2/Y}.}{c8:fig:statediagram}

\noindent Each state bubble shows the current state and the output as \code{Q1Q2/Y}, and each arrow
is labelled with the input \code{X} that triggers it. Each state loops back to itself while
\code{X=0}, and a synchronized button pulse (\code{X=1}) advances to the next state, wrapping from
\code{STATE_3} back to \code{STATE_0}.

\needspace{14\baselineskip}
\subsection{State table}
\label{c8:table}
Every combination of the current state and input, giving the next state \code{\{Q1+, Q2+\}} and
output \code{Y}:

\begin{narrowtable}{@{}cccccc@{}}
  \thead{Q1} & \thead{Q2} & \thead{X} & \thead{Q1+} & \thead{Q2+} & \thead{Y} \\
  \midrule
  0 & 0 & 0 & 0 & 0 & 0 \\
  0 & 0 & 1 & 0 & 1 & 0 \\
  0 & 1 & 0 & 0 & 1 & 0 \\
  0 & 1 & 1 & 1 & 1 & 0 \\
  1 & 0 & 0 & 1 & 0 & 1 \\
  1 & 0 & 1 & 0 & 0 & 1 \\
  1 & 1 & 0 & 1 & 1 & 0 \\
  1 & 1 & 1 & 1 & 0 & 0 \\
\end{narrowtable}

\subsection{Karnaugh-derived logic}
\label{c8:kmap}
Treating \code{Q1+} and \code{Q2+} each as a Boolean function of \code{\{Q1, Q2, X\}} and running
them through a Karnaugh map, the same technique as Chapter~\ref{c2:ch}:

\begin{console}
Q1+ = Q1X' + Q2X
Q2+ = Q2X' + Q1'X
\end{console}

\noindent Notice the shape of both:
\begin{itemize}
  \item each is ``hold when \code{X} is \code{0}, take something from the other bit when \code{X} is
    \code{1}.''
  \item the hold halves, \code{Q1X'} and \code{Q2X'}, really are mirror images of each other.
  \item the advance halves are \code{Q2X} and \code{Q1'X}, and that complement on the second one is
    not a typo: it is what makes the sequence a Gray code rather than a plain rotation.
\end{itemize}
The output \code{Y} reduces further, straight from the state table, with no dependency on \code{X}
at all, confirming this is a valid Moore machine:

\begin{console}
Y = Q1Q2'
\end{console}

\subsection{Realizing it as a gate network}
\label{c8:gates}
\code{\{Q1, Q2\}} are each a D flip-flop, exactly like Chapter~\ref{c3:ch}'s register:
\code{Q1+}/\code{Q2+} feed the D inputs and \code{Q1}/\code{Q2} are read back from the Q outputs.
\code{Y} is a single two-input AND gate with one input inverted, reading the current state. That is
the entire machine: two flip-flops and a handful of gates computing \code{Q1+}, \code{Q2+} and
\code{Y} from \code{\{Q1, Q2, X\}}.

\section{From Gate Network to CircuitVerse}
\label{c8:a3}
Before writing any of this in VHDL, realize it by hand in CircuitVerse
(\url{https://circuitverse.org/simulator}), exactly as every gate network in Chapter~\ref{c1:ch}
was: double-flop synchronize \code{button_n} and \code{reset_n}, edge-detect the synchronized button
to produce \code{X}, wire \code{Q1+}/\code{Q2+} into a pair of D flip-flops, and wire \code{Y}
straight into the LED.

\bookimage{1.0}{lectures/L08/appendix/images/fsm_gate_network.png}{The four-state Moore machine as
a gate network in CircuitVerse, with the double-flop synchronization and edge detection of
\code{button_n} on the left and \code{Y} driving the LED on the right.}{c8:fig:gatenetwork}

\noindent Reading it left to right:
\begin{itemize}
  \item \code{button_n} passes through three synchronizer flip-flops
    (\code{button_s1_n}/\allowbreak\code{button_s2_n}/\allowbreak\code{button_s3_n}).
  \item An AND gate derives \code{button_edge_s2} (this circuit's \code{X}) from the second and
    third stages.
  \item The four vertical rails labelled \code{Q1Q2XX'} carry the state bits and the edge pulse
    across to the gates on the right, which implement \secref{c8:a2}'s \code{Q1+}/\code{Q2+}
    equations.
  \item The two labelled flip-flops store \code{Q1}/\code{Q2} and feed the final AND/NOT gates
    producing \code{Y}, wired to \code{led}.
  \item \code{reset_n} is synchronized by its own pair of flip-flops at bottom left, exactly like
    every other synchronized reset in this book.
\end{itemize}
Open the live circuit yourself: \file{lectures/L08/appendix/images/fsm_gate_network.cv}. Step it by
hand, one button press at a time, and confirm it visits \code{STATE_0 -> STATE_1 -> STATE_2 ->
STATE_3 -> STATE_0} in order, with the LED lighting only on \code{STATE_3}.
